\documentclass[10pt]{article}
\usepackage{icmlko}

\icmlmeta{IP 파운데이션 월드모델}{504}{배경이 바뀌고 있었다 --- 스텝 501 의 마지막 문장이 선 근거는 배경 풀의 성질이다}{그리고 그 자리의 부호는 결론이 아니라 곡선이다: 배경과 구성을 둘 다 고정해도 $n$ 을 따라 넘어가고, 진짜 군집을 대면 관측 구간 스무 칸이 전부 「못 가른다」다}

\begin{document}

\twocolumn[
\icmlheader

\begin{abstract}
\textbf{이것은 스텝 501 의 다섯째 주장에 대한 정정본이다.} 501 은 「층 ③의 정보는
도메인 이름표로 \emph{대체 가능}하다」를 \emph{살아남는 문장}으로 실었고, 근거는
464행 표에서 「도메인 지시자 위에 층 ③을 얹으면 $\Delta\rho=-$0.0075 로 \textbf{빠진다}」
하나였다. \textbf{그 자리에서 표만 바뀐 게 아니라 배경 풀도 같이 바뀌고 있었다.}
층 ③은 \emph{도메인 배경 대비 편차}라서 배경 풀을 갈면 그 열의 모든 원소가 바뀐다.
\textbf{행 464개 · 지시자 9열 · 도메인 구성 · 재표집 설계를 전부 같게 두고 배경 풀만
813개체에서 3,896개체로 바꾸면} 그 수가 $-$0.007687(문턱 0.006557, \textbf{「뺀다」},
여유 1.172배)에서 \textbf{$-$0.005928}(문턱 0.007622, \textbf{「못 가른다」}, 여유
0.778배)로 내려앉는다. \textbf{바뀐 것은 그 열의 해시 하나다.} 더 나아가 우리는 그
부호를 \textbf{표본 크기 $n$ 의 함수}로 사전등록하고 쟀다 --- 배경을 3,896에 못 박고
도메인 구성까지 고정한 채 부분표집 씨앗 30개로 두 팔을 올리면, 부호는 두 팔 모두에서
음수에서 양수로 \textbf{매끄럽지 않게} 넘어간다(영교차 $n^{*}=$ 561.3 과 377.8).
\textbf{그리고 개체가 아니라 도메인을 군집으로 놓으면 관측 구간 스무 칸이 전부 「못
가른다」다 --- 예외가 없다.} 포화 적합의 점근값조차 그 문턱의 39.1\,\%다.
\textbf{그러므로 이 자로는 「층 ③이 도메인 이름표의 대리변수인가」를 물을 수 없다}는
것이 이 스텝의 결론이고, 501의 다섯째 주장과 그것을 뒤집었다던 문장은 \textbf{둘 다}
「못 가른다」로 내린다.
\end{abstract}
]

\section{무엇이 틀렸나}

501 의 다섯째 주장은 \textbf{하나의 수}에 걸려 있었다.

\begin{itemize}
\item[(가)] 도메인 지시자만 넣으면 $+$0.087630 을 번다 --- \textbf{그대로다.}
\item[(나)] 그 \emph{위에서} 층 ③은 $-$0.0075 로 \textbf{빠진다} --- \textbf{여기가 문제다.}
\end{itemize}

(나)의 근거가 된 464행 표는 \emph{배경 풀 813개체} 위에서 지어졌다. 그 뒤 표를 2,050행으로
키운 판정은 \emph{배경 풀 3,896개체} 위에서 지어졌다. 두 판정을 「자를 고정하고 표만
바꿨다」로 이으면 \textbf{배경이 조용히 같이 움직인다}. 층 ③은 도메인 배경 대비 편차라
배경 풀을 갈면 \textbf{그 열 전체가 다른 열이 된다}.

\section{배경만 갈았다}

행·열·도메인 구성·재표집 설계(4,000회 · 같은 씨앗)를 전부 같게 두고 배경 풀만 바꿨다.

\begin{center}
\begin{tabular}{lrrl}
\hline
배경 풀 & $\Delta\rho$ & 문턱 $T$ & 판정 \\
\hline
813 개체  & $-$0.007687 & 0.006557 & 뺀다 \\
3,896 개체 & $-$0.005928 & 0.007622 & \textbf{못 가른다} \\
\hline
\end{tabular}
\end{center}

행 수 464 = 464, 지시자 9열 = 9열, 도메인별 행 수 동일, 재표집 동일. 두 팔에서 다른
것은 \textbf{그 편차 열의 해시 하나}다. 여유는 1.172배에서 \textbf{0.778배}로 내려간다.
$-$0.007687 은 앞선 판정의 수를 자릿수까지 재현한 값이므로, 이 차이는 재현 실패가
아니라 \textbf{배경 풀의 효과}다.

\section{그리고 그 부호는 곡선이다}

배경 풀을 3,896에 못 박고, 도메인 구성까지 팔마다 고정하고, 부분표집 씨앗 30개로
$n$ 을 올렸다. 두 팔은 구성도 지시자 열 수도 다르므로 \textbf{절대값을 한 문장에 잇지
않는다}.

\begin{center}
\begin{tabular}{lrrr}
\hline
팔 & $n$ 범위 & 양끝 $m(n)$ & 영교차 $n^{*}$ \\
\hline
㉠ (9열)  & 302--851   & $-$0.003767 $\to$ $+$0.002300 & 561.3 \\
㉡ (10열) & 301--2,050 & $-$0.002958 $\to$ $+$0.007112 & 377.8 \\
\hline
\end{tabular}
\end{center}

부호는 두 팔 모두에서 넘어간다. \textbf{배경도 구성도 안 움직였는데} 넘어간다. 그러므로
한쪽 끝의 값을 「대체 가능하다」로, 다른 쪽 끝의 값을 「대체 가능하지 않다」로 읽는 것은
\textbf{같은 곡선의 두 점을 두 개의 결론으로 읽은 것}이다.

\section{그런데 이 곡선은 판정할 수 없다}

개체를 군집으로 쓴 재표집은 개체당 한 행이라 사실상 독립 행 재표집이다. \textbf{도메인을
군집으로 놓으면} 두 팔 스무 칸이 \textbf{전부} 「못 가른다」가 된다 --- 예외 0칸. 그리고
개체 군집 아래서도 판정은 자료를 안 따라간다: $m(n)$ 이 $n=$ 1,001에서 2,050까지
$+$0.006065에서 $+$0.007112로 거의 평평한 동안 판정은 「든다 · 못 가른다 · 못 가른다 ·
든다 · 든다 · 든다 · 든다」로 튄다. \textbf{튄 것은 $\Delta\rho$ 가 아니라 문턱이다}
(0.004639 · 0.012176 · 0.008217 · 0.003852 $\dots$). \textbf{판정을 정한 것은 자료가
아니라 재표집의 잡음이다.}

곡선이 평평해지는 $n$ 이 있는지도 문턱까지 미리 등록해 쟀다(한 $e$배마다 $\Delta\rho$
변화가 카드 문턱 0.00353보다 작고, 동시에 그 기울기의 95\,\% 구간이 0을 품을 것).
\textbf{답이 팔마다 갈렸다} --- ㉠은 652부터 평평하지만 \emph{그 뒤 칸이 둘뿐}이고,
㉡은 끝까지 평평해지지 않는다. 포화 적합 $m(n)=a-b\,n^{-1/2}$ 의 점근값은 ㉡에서
$+$0.013703 $[+$0.012252, $+$0.015114$]$ 인데, 이는 도메인 군집 문턱 0.035077의
\textbf{39.1\,\%}다. \textbf{곡선을 $n\to\infty$ 로 늘려 읽어도 이 자는 답에 못 닿는다.}

\section{자는 넷이 아니라 하나였다}

같은 판정을 네 자로 확인했다고 적어 온 자리를 열어 보면, 자들 사이에 \textbf{논리적
포함관계}가 있다. 순열 영분포의 통과 문턱은 $-$0.000493 --- 카드 문턱의 $-$0.1397배로
사실상 0이라 그 자는 「$\Delta\rho>0$」과 같고, 합산 자의 문턱은 0.020060이다. 합산이
참이면 순열은 자동으로 참이다. 부 표적 자는 부호만 보므로 문턱이 0이고 역시 포함된다.
\textbf{독립된 자의 수는 1이다.} 「자 넷이 다 같은 답을 냈다」는 \textbf{하나의 답을 넷으로
센 것}이다.

\section{채택 규칙을 3값으로 고쳤다}

채택 규칙의 네 번째 조건(더 싼 대체물이 있나)은 두 하위검사를 참/거짓 둘로만 냈다.
그래서 하위검사가 \emph{문턱 띠 안}(= 못 가른다)일 때 조용히 「대체물이 안 든다」와
「그 위에서 층이 죽는다」로 접혔고, \textbf{두 하위검사가 둘 다 「못 가른다」인데 「채택」이
나오는} 자리가 있었다. 참/거짓/\textbf{모른다} 셋으로 바꾸고 클레이니 논리곱으로 이으니
그 자리가 \textbf{보류}로 바뀐다. 「모른다」는 「통과」가 아니다.

\section{그래서 무엇이 남나}

\begin{itemize}
\item 501의 다섯째 주장 「③의 정보는 도메인 이름표로 대체 가능하다」 --- \textbf{「못 가른다」로 내린다.}
      그 근거 수는 배경 풀 813의 성질이었다.
\item 그것을 뒤집었다던 반대 문장(아직 병합되지 않은 가지에만 있다) --- \textbf{같이 「못 가른다」로 내린다.}
      근거의 여유가 1.172배와 1.166배로 \textbf{같은 자리}였다.
\item \textbf{남는 문장은 하나다}: 이 자로는 그 물음을 물을 수 없다.
\end{itemize}

\section{반증조건}

\begin{itemize}
\item 도메인을 군집으로 놓은 재표집에서 어떤 $n$ 이든 $|\Delta\rho|$ 가 그 문턱을 넘으면
      이 스텝의 결론(「물을 수 없다」)이 무너진다. 관측 구간에서 그 여유의 최대값은
      \textbf{0.39배}이므로, 문턱을 넘으려면 지금의 2.6배가 필요하다.
\item 배경 풀을 키우면서 \emph{내용}을 고정하는 설계(같은 도메인 구성으로 배경만 증량)에서
      위 표의 차이가 사라지면, 「배경 풀의 효과」라는 이 스텝의 진단이 무너진다.
      \textbf{이 스텝은 배경의 크기와 내용을 못 갈랐다.}
\item $n>$ 2,050 의 표에서 곡선이 실제로 평평해지고 그 평평한 값이 도메인 군집 문턱을
      넘으면, 「물을 수 없다」는 「그때는 물을 수 있다」로 좁혀진다.
\end{itemize}

\section{이 스텝이 스스로 신고하는 것}

배경 풀의 \emph{크기}와 \emph{내용}을 못 갈랐다(813$\to$3,896은 개체 수도 구성도 같이
바꾼다). 두 팔의 높이 차를 「구성 효과」로 못 읽는다(지시자 열이 9와 10으로 다르다).
재표집은 각 칸의 씨앗 0 표본 하나 위에서만 돌았으므로 문턱 자체의 씨앗 변동은 안 쟀다.
배경만 갈아 낸 표는 \textbf{사전등록에 없는 사후 측정}이다 --- 곡선의 닻이 앞선 판정의
수와 어긋나서 원인을 못 박으려고 뒤에 돌렸고, 판정은 거기서 세우지 않았다.

\end{document}
